\chapter{Performance Metrics and the PRAM Model}

% -------------------------------------------------------------------
\section*{Learning Outcomes}
% -------------------------------------------------------------------
By the end of this chapter, you will be able to:
\begin{itemize}
  \item Define speedup, cost, efficiency, and parallelism overhead, and
        compute them for a given parallel algorithm.
  \item Distinguish between strong scaling and weak scaling, and explain
        when each is the appropriate measure.
  \item Apply Amdahl's Law to derive an upper bound on speedup for a fixed
        problem size, and explain why that bound is pessimistic.
  \item Apply Gustafson's Law to compute speedup under a scaled-workload
        model, and contrast it with Amdahl's Law.
  \item Derive the isoefficiency function for a parallel system and use it
        to classify the system as strongly or weakly scalable.
  \item Describe the PRAM model of computation and its three synchronisation
        variants (EREW, CREW, CRCW).
  \item Analyse a simple PRAM algorithm (parallel addition) for time, work,
        cost, speedup, and efficiency.
\end{itemize}

% -------------------------------------------------------------------
\section{Revisiting Performance Metrics}
% -------------------------------------------------------------------

Understanding how a parallel program performs requires more than just measuring
wall-clock time. We need a vocabulary that captures \emph{how efficiently}
resources are used and \emph{how well} a solution scales as resources and
problem sizes change.

\begin{ppdefn}
Let $T_1$ denote the time taken on a single sequential processor, and $T_p$ the
time taken on a parallel system with $p$ processors. The \ppterm{speedup} is
\[
  S_p \;:=\; \frac{T_1}{T_p}.
\]
Ideal speedup on $p$ processors is $S_p = p$.
\end{ppdefn}

\begin{ppdefn}
The \ppterm{cost} of a parallel execution is
\[
  C_p \;:=\; p \times T_p.
\]
It captures the total processor-time expended. A parallel algorithm is
\ppterm{cost-optimal} when $p \cdot T_p = O(T_1)$, i.e., the total work
done is asymptotically no more than the sequential baseline.
\end{ppdefn}

\begin{ppexample}
Suppose $T_1 = 100$ and $T_8 = 20$. Then $S_8 = 5$ and $C_8 = 160$. The
parallel system used 60\% more total processor-time than a sequential run,
so it is \emph{not} cost-optimal.
\end{ppexample}

\subsection{Overhead and Efficiency}

Running on $p$ processors introduces overhead from communication,
synchronisation, and idle time. These reduce efficiency below the ideal.

\begin{ppdefn}
The \ppterm{parallelism overhead} $T_o$ (also written $\bar{o}_p$) is the
extra work introduced by running in parallel:
\[
  T_o \;:=\; C_p - T_1 \;=\; p \cdot T_p - T_1.
\]
\end{ppdefn}

\begin{ppdefn}
\ppterm{Efficiency} measures how effectively processors are utilised:
\[
  E_p \;:=\; \frac{T_1}{C_p} \;=\; \frac{S_p}{p}.
\]
Efficiency is 1 when $S_p = p$ (ideal), and falls as overhead grows.
\end{ppdefn}

\begin{ppexample}
Consider two configurations: $S_8 = 7$ and $S_{64} = 20$.
\[
  E_8 = \frac{7}{8} = 0.875, \qquad E_{64} = \frac{20}{64} \approx 0.31.
\]
The 64-processor run achieves higher absolute speedup but is far less efficient:
nearly 70\% of processor time is wasted on overhead.
\end{ppexample}

\begin{pppitfall}
Raw speedup numbers can be misleading. Always report efficiency alongside
speedup---a speedup of 20 on 64 processors is underwhelming, not impressive.
\end{pppitfall}

% -------------------------------------------------------------------
\section{Scalability}
% -------------------------------------------------------------------

Raw speedup is a snapshot at one problem size and one processor count.
\ppterm{Scalability} asks how performance evolves as \emph{both} problem
size $n$ and processor count $p$ change.

Formally, we now treat time as $T(n, p)$---a function of both parameters---and
interpret speedup, efficiency, and cost accordingly.

\subsection{Strong Scaling}

\ppterm{Strong scaling} studies what happens when the problem size $n$ is
held fixed and $p$ increases.

\begin{ppdefn}
Ideal strong scaling: $T(n, p) = T(n, 1)/p$ for fixed $n$.
\end{ppdefn}

In practice the actual parallel time decomposes as
\[
  T(n, p) = T_{\text{comp}} + T_{\text{comm}} + T_{\text{sync}} + T_{\text{idle}}.
\]
Only $T_{\text{comp}}$ shrinks with $p$; the remaining terms are overheads that
grow with $p$ and eventually dominate, causing strong scaling to break down.

\subsection{Weak Scaling}

\ppterm{Weak scaling} studies what happens when each processor is assigned a
fixed workload $m$, so the total problem size grows as $mp$ when $p$ increases.

\begin{ppdefn}
Ideal weak scaling: $T(mp, p) = T(m, 1)$---runtime stays constant as both
problem size and processor count grow proportionally.
\end{ppdefn}

\begin{ppexample}
A weakly scaling system might show:
\begin{center}
\begin{tabular}{ccc}
\toprule
$p$ & Problem size & Time (s) \\
\midrule
1  & 1 million & 10 \\
2  & 2 million & 11 \\
4  & 4 million & 13 \\
\bottomrule
\end{tabular}
\end{center}
Runtime grows only slightly---the small increase is due to communication
overhead that scales sub-linearly.
\end{ppexample}

% -------------------------------------------------------------------
\section{Amdahl's Law}
% -------------------------------------------------------------------

\begin{ppkey}
Amdahl's Law bounds the speedup achievable for a fixed problem size when
only a fraction of the computation is parallelisable.
\end{ppkey}

Suppose a fraction $f$ of the total work is inherently sequential and
$(1 - f)$ is fully parallelisable. Then:
\[
  T_p = T_1\!\left(f + \frac{1-f}{p}\right),
  \qquad
  S_p = \frac{1}{f + (1-f)/p},
  \qquad
  S_\infty = \frac{1}{f}.
\]

\begin{ppexample}
If only 10\% of the program is sequential ($f = 0.1$), the maximum possible
speedup is $S_\infty = 1/0.1 = 10$, regardless of how many processors are
added.
\end{ppexample}

\begin{figure}[h]
\centering
\begin{tikzpicture}
\begin{axis}[ppplot,
  xlabel={Number of processors $p$},
  ylabel={Speedup $S_p$},
  xmin=1, xmax=64, ymin=1, ymax=22,
  xtick={1,8,16,32,64},
]
\addplot[black, dashed, thick, domain=1:22, samples=22] {x};
\addlegendentry{Ideal ($S_p = p$)}
\addplot[color=PPBlue,   thick, domain=1:64, samples=100] {1/(0.05 + 0.95/x)};
\addlegendentry{$f = 0.05$, $S_\infty = 20$}
\addplot[color=PPGood,   thick, domain=1:64, samples=100] {1/(0.1  + 0.9/x)};
\addlegendentry{$f = 0.10$, $S_\infty = 10$}
\addplot[color=orange,   thick, domain=1:64, samples=100] {1/(0.25 + 0.75/x)};
\addlegendentry{$f = 0.25$, $S_\infty = 4$}
\addplot[color=PPWarn,   thick, domain=1:64, samples=100] {1/(0.5  + 0.5/x)};
\addlegendentry{$f = 0.50$, $S_\infty = 2$}
\end{axis}
\end{tikzpicture}
\caption{Amdahl's Law: speedup saturates at $1/f$ regardless of processor count.
         Even with only 5\% sequential code the ceiling is~20.}
\label{fig:amdahl-speedup}
\end{figure}

\begin{pppitfall}
Improving the parallel portion of a program yields diminishing returns if the
sequential fraction remains significant. Amdahl's Law is the formal statement
of this limit. Optimising the hot loop by $10\times$ is irrelevant if the
sequential setup and teardown dominate.
\end{pppitfall}

Amdahl's Law applies to fixed problem sizes. It does not capture what happens
when workload grows with the number of processors---that is Gustafson's domain.

% -------------------------------------------------------------------
\section{Gustafson's Law}
% -------------------------------------------------------------------

Amdahl's Law is pessimistic because it holds problem size fixed. In practice,
when more processors become available, practitioners scale up the problem size
to do more useful work.

\begin{ppkey}
As $p$ increases, the opportunity for parallelisation also increases. Speedup
can grow with $p$ if workload is scaled proportionally.
\end{ppkey}

Let $f \cdot T(n, p)$ be the time spent in the sequential portion on $p$
processors, and $(1 - f) \cdot T(n, p)$ the time spent in the parallel
portion. The sequential time to solve the same scaled problem is
\[
  T(n, 1) = f \cdot T(n, p) + p \cdot (1 - f) \cdot T(n, p).
\]
Therefore:
\begin{ppdefn}
\ppterm{Gustafson's speedup}:
\[
  S_p = \frac{T(n, 1)}{T(n, p)} = f + (1 - f) \cdot p.
\]
\end{ppdefn}

Unlike Amdahl's Law, this speedup grows linearly with $p$ as long as $f$ is
small---and $f$ shrinks as the parallelisable portion of the workload grows.

\begin{figure}[h]
\centering
\begin{tikzpicture}
\begin{axis}[ppplot,
  xlabel={Number of processors $p$},
  ylabel={Speedup $S_p$},
  xmin=1, xmax=16, ymin=1, ymax=16,
  xtick={1,4,8,12,16},
]
\addplot[black, dashed, thick, domain=1:16, samples=16] {x};
\addlegendentry{Ideal}
\addplot[color=PPWarn, thick, domain=1:16, samples=100] {1/(0.1 + 0.9/x)};
\addlegendentry{Amdahl, $f=0.1$}
\addplot[color=PPBlue, thick, domain=1:16, samples=16]  {0.1 + 0.9*x};
\addlegendentry{Gustafson, $f=0.1$}
\addplot[color=PPGood, thick, domain=1:16, samples=100] {1/(0.25 + 0.75/x)};
\addlegendentry{Amdahl, $f=0.25$}
\addplot[color=orange, thick, domain=1:16, samples=16]  {0.25 + 0.75*x};
\addlegendentry{Gustafson, $f=0.25$}
\end{axis}
\end{tikzpicture}
\caption{Amdahl vs.\ Gustafson speedup for the same serial fraction $f$.
         Amdahl (fixed workload) saturates; Gustafson (scaled workload) grows
         linearly---the key insight is that scaling the problem size reveals
         more parallelism.}
\label{fig:amdahl-vs-gustafson}
\end{figure}

% -------------------------------------------------------------------
\section{IsoEfficiency}
% -------------------------------------------------------------------

Gustafson's Law tells us speedup can grow, but does not say \emph{by how much
the workload must grow} to keep efficiency constant as $p$ increases. The
\ppterm{isoefficiency function} answers this question.

\begin{ppdefn}
The \ppterm{isoefficiency function} gives the rate at which the workload $W$
(equivalently $T_1$) must grow as a function of $p$ to maintain a constant
efficiency $E_c$:
\[
  E(n, p) = \frac{T_1}{p \cdot T(n, p)} = E_c.
\]
\end{ppdefn}

Reformulating using $T_o = p \cdot T_p - T_1$:
\[
  E = \frac{T_1}{T_1 + T_o} \implies W = \frac{E_c}{1 - E_c} \cdot T_o.
\]

This gives the minimum workload required to maintain efficiency $E_c$ given
the overhead $T_o$. If $T_o$ grows faster than $W$, efficiency falls; the
workload must grow at least as fast as $T_o$ to hold efficiency constant.

\begin{ppnote}
If $n$ need not grow with $p$ to maintain constant efficiency, the system is
\ppterm{strongly scalable}. If $n$ must grow with $p$, the system is
\ppterm{weakly scalable}. The isoefficiency function quantifies the rate of
required growth.
\end{ppnote}

\begin{ppexample}
Suppose $T_o = p \log p$ and we want to maintain $E_c = 0.8$. Then:
\[
  W = \frac{0.8}{0.2} \cdot T_o = 4 T_o = \Theta(T_o) = \Theta(p \log p).
\]
The workload must grow at least as $p \log p$ to keep efficiency at 80\%.
If the workload grows slower, efficiency falls; if faster, efficiency may
improve.
\end{ppexample}

% -------------------------------------------------------------------
\section{A Running Example: Matrix Multiplication}
% -------------------------------------------------------------------

We tie all the metrics together with a concrete parallel matrix multiplication
to build intuition.

Assume $n \times n$ matrices. Sequential time: $T_1(n) = n^3$. Parallel time:
\[
  T_p = \frac{n^3}{p} + \underbrace{(n + 1000 \log p)}_{\text{overhead}}.
\]
The overhead captures the sequential setup (reading input, $O(n)$) and
communication costs ($O(\log p)$ rounds at constant bandwidth).

\subsection{Speedup and Efficiency at $n = 1000$}

For $n = 1000$, computation dominates: $n^3 = 10^9 \gg n + 1000 \log p$.
Therefore $T_p \approx n^3/p$, giving $S_p \approx p$ and $E_p \approx 1$.

\subsection{Applying Amdahl's Law}

The serial fraction is $f = n / n^3 = 1/n^2$. For $n = 1000$, $f = 10^{-6}$.
Including the communication overhead for a fixed $p = 64$:
\[
  n + 1000 \log 64 = 1000 + 6000 = 7000 \implies f_{\text{eff}} \approx 7 \times 10^{-6}.
\]
So $S_{64} \approx 1/f_{\text{eff}} \approx 64$---ideal. But if $n = 100$:
$S_{64} \approx 46$ and $E_{64} \approx 0.70$, showing that smaller problem
sizes expose overheads.

\subsection{Strong and Weak Scaling}

For $n = 1000$, strong scaling is excellent up to $p = 64$ since $T_p \approx T_1/p$.
Strong scaling degrades for $n = 100$ because the overhead terms become comparable
to computation.

For weak scaling, fix workload per processor at $W = 10^6$ operations:

\begin{center}
\begin{tabular}{ccc}
\toprule
$p$ & $n$ (from $n^3 = p \cdot 10^6$) & $T_p$ (approx) \\
\midrule
1  & 100 & $10^6$ \\
8  & 200 & $10^6 + 3200$ \\
64 & 400 & $10^6 + 6400$ \\
\bottomrule
\end{tabular}
\end{center}

Runtime stays approximately constant in $p$---the small increase comes only from
the $1000 \log p$ communication term. This is very good weak scaling.

\subsection{Gustafson's Law Applied}

For $p = 64$, $n = 400$: the serial and overhead portion is $400 + 6000 = 6400$.
\[
  f = \frac{6400}{1006400} \approx 0.0064,
  \qquad
  S_{64} = f + (1-f) \cdot p = 0.0064 + 0.9936 \times 64 \approx 63.6.
\]
The serial fraction is negligibly small, so speedup remains near-ideal even
under Gustafson's analysis.

\subsection{IsoEfficiency}

The overhead is $T_o = np + 1000p \log p$. For large $n$, the term $np$
dominates. Setting $W = n^3 = \Theta(T_o) = \Theta(np)$:
\[
  n^3 = \Theta(np) \implies n^2 = \Theta(p) \implies n = \Theta(\sqrt{p})
  \implies W = \Theta(p^{3/2}).
\]
To maintain constant efficiency, the workload must grow as $p^{3/2}$. This is
acceptable but not ideal (ideal would be $\Theta(p)$); for every additional
processor, the problem size must grow super-linearly.

\begin{pppitfall}
An isoefficiency of $\Theta(p^2)$ or worse means the system scales poorly:
problem size must grow quadratically just to keep efficiency from falling.
Always compute the isoefficiency function before declaring a parallel
algorithm ``efficient.''
\end{pppitfall}

% -------------------------------------------------------------------
\section{A Note on Superlinear Speedup}
% -------------------------------------------------------------------

Theoretically, speedup cannot exceed $p$: the ideal parallel time is $T_1/p$.
In practice, however, \ppterm{superlinear speedup} ($S_p > p$) is observed.

\begin{ppnote}
Superlinear speedup arises when decomposing the problem across processors
exposes \emph{architectural} benefits not captured by asymptotic analysis.
The most common cause is cache effects: a sub-problem assigned to one
processor fits entirely in its L2 or L3 cache, whereas the full problem did
not fit in any single processor's cache. Effective memory latency drops from
DRAM latency to cache latency, accelerating computation beyond what additional
cores alone would predict. Smart data decomposition---both algorithmic and
hardware-aware---is the mechanism.
\end{ppnote}

% -------------------------------------------------------------------
\section{The PRAM Model of Computation}
% -------------------------------------------------------------------

To reason formally about parallel algorithms, we need a clean theoretical
model that abstracts away architectural details. The \ppterm{PRAM}
(Parallel Random Access Machine) model is the standard choice.

\begin{ppdefn}
A \ppterm{PRAM} consists of:
\begin{itemize}
  \item An arbitrary number of synchronous processors, each with a unique
        integer identifier.
  \item Unbounded local memory per processor.
  \item A shared global memory accessible by any processor in a single step
        (constant-time access---realistic for asymptotic analysis, unrealistic
        in practice).
\end{itemize}
Processors execute in lock-step. Performance is measured by:
\begin{itemize}
  \item \textbf{Time}: steps taken by the longest-running processor.
  \item \textbf{Work}: total operations across all processors.
  \item \textbf{Cost}: $p \times \text{Time}$.
\end{itemize}
\end{ppdefn}

\begin{ppnote}
The PRAM's constant-time shared memory access is unrealistic but useful: it
lets us focus on algorithmic parallelism rather than memory-hierarchy effects.
Real systems pay variable latency for remote memory; the PRAM model establishes
theoretical upper bounds that guide design even if they cannot be achieved
exactly.
\end{ppnote}

\subsection{Concurrent Access: The Synchronisation Model}

Multiple processors reading or writing the same shared location simultaneously
raises consistency questions. The PRAM model is parameterised by its
\ppterm{synchronisation model}, which specifies the rules for concurrent access.

\begin{ppdefn}
\ppterm{Exclusive Read, Exclusive Write (EREW)}: at any step, at most one
processor may read and at most one may write any given shared location.
Concurrent accesses to the same location are forbidden.
\end{ppdefn}

EREW is the most restrictive model. Algorithms designed for EREW are safe on
essentially any hardware, but they may need extra steps to route data through
intermediary locations when multiple processors need the same value.

\begin{ppdefn}
\ppterm{Concurrent Read, Exclusive Write (CREW)}: multiple processors may
simultaneously read a shared location, but writes remain exclusive.
\end{ppdefn}

CREW matches the common producer-consumer pattern: one writer updates a
shared result while many readers consume it. Most PRAM algorithms studied in
practice assume CREW.

\begin{ppdefn}
\ppterm{Concurrent Read, Concurrent Write (CRCW)}: multiple processors may
read and write the same location simultaneously. A \ppterm{write conflict
resolution rule} determines the outcome.
\end{ppdefn}

The three standard write conflict rules for CRCW are:

\begin{description}
  \item[Priority CW.] The processor with the highest priority (typically the
    lowest index) wins; its value is written.
  \item[Common CW.] The write succeeds only if all concurrent writers attempt
    to write the \emph{same} value. If any writer disagrees, the write fails
    and the location retains its old value.
  \item[Arbitrary/Random CW.] One of the conflicting writes is chosen
    arbitrarily (or randomly); the rest are discarded.
\end{description}

\begin{ppnote}
The models form a hierarchy by computational power per step (increasing left
to right):
\[
  \text{EREW} \;\leq\; \text{CREW} \;\leq\; \text{Common CRCW}
  \;\leq\; \text{Arbitrary CRCW} \;\leq\; \text{Priority CRCW}.
\]
A stronger model can simulate a weaker one with no overhead, but not vice
versa. An algorithm correct under EREW is correct under all stronger models;
an algorithm relying on Priority CRCW may not port to EREW without extra
steps. (Diagram courtesy: S.~Khuller.)
\end{ppnote}

\begin{ppexample}
Common CW arises naturally in \emph{leader election}: every processor writes
its agreed-upon leader to a shared register. If all agree, the write succeeds;
a failed write reveals that consensus was not reached.
\end{ppexample}

% -------------------------------------------------------------------
\section{Parallel Addition on the PRAM}
% -------------------------------------------------------------------

As a first algorithm on the PRAM, consider summing $n$ elements of an array
$A[0 \ldots n-1]$ under the CREW model.

\textbf{Algorithm.} Use a binary-tree reduction
(Figure~\ref{fig:pram-tree} shows an example for $n=8$):
\begin{enumerate}
  \item Phase~1 ($p = n$ processors): each processor $i$ copies $A[i]$ into
        a working buffer $B[i]$.
  \item Repeat: halve $p$. Each active processor $i < p$ computes
        $B[i] \leftarrow B[2i] + B[2i+1]$.
  \item After $\lceil \log_2 n \rceil$ rounds, $B[0]$ holds the sum.
\end{enumerate}

\begin{figure}[h]
\centering
\begin{tikzpicture}[
  >=Stealth,
  box/.style={draw=PPBlue!60, fill=PPLight, rounded corners=2pt,
              minimum width=0.85cm, minimum height=0.42cm, font=\small},
  sum/.style={draw=PPGood!80, fill=PPGood!12, rounded corners=2pt,
              minimum width=0.95cm, minimum height=0.42cm,
              font=\small\bfseries},
  lbl/.style={font=\footnotesize\itshape, text=PPGray},
]
% Row 0 — initial values (8 processors)
\node[box] (a0) at (0*1.5, 0)   {3};
\node[box] (a1) at (1*1.5, 0)   {7};
\node[box] (a2) at (2*1.5, 0)   {11};
\node[box] (a3) at (3*1.5, 0)   {15};
\node[box] (a4) at (4*1.5, 0)   {2};
\node[box] (a5) at (5*1.5, 0)   {4};
\node[box] (a6) at (6*1.5, 0)   {8};
\node[box] (a7) at (7*1.5, 0)   {6};
% Row 1 — round 1 (4 processors)
\node[sum] (b0) at (0.75,  1.3) {10};
\node[sum] (b1) at (3.75,  1.3) {26};
\node[sum] (b2) at (6.75,  1.3) {6};
\node[sum] (b3) at (9.75,  1.3) {14};
% Row 2 — round 2 (2 processors)
\node[sum] (c0) at (2.25,  2.6) {36};
\node[sum] (c1) at (8.25,  2.6) {20};
% Row 3 — round 3 (1 processor)
\node[sum] (d0) at (5.25,  3.9) {56};
% Edges: row 0 -> row 1
\foreach \src/\dst in {a0/b0, a1/b0, a2/b1, a3/b1, a4/b2, a5/b2, a6/b3, a7/b3}
  \draw[->] (\src) -- (\dst);
% Edges: row 1 -> row 2
\foreach \src/\dst in {b0/c0, b1/c0, b2/c1, b3/c1}
  \draw[->] (\src) -- (\dst);
% Edges: row 2 -> row 3
\draw[->] (c0) -- (d0);
\draw[->] (c1) -- (d0);
% Row labels
\node[lbl, left=0.25cm of a0] {Round 0};
\node[lbl, left=0.25cm of b0] {Round 1};
\node[lbl, left=0.25cm of c0] {Round 2};
\node[lbl, left=0.25cm of d0] {Round 3};
\end{tikzpicture}
\caption{Binary-tree reduction for $n=8$ elements. Each round halves the active
         processor count; the sum~56 is reached in $\lceil\log_2 8\rceil = 3$
         rounds. Processors owning a right-child value (\emph{odd} indices) are
         idle after Round~1.}
\label{fig:pram-tree}
\end{figure}

\begin{ppprogram}{PRAM parallel sum (binary-tree reduction)}{prog:pram-sum}
\begin{lstlisting}[language={}]
/* Pseudocode: n processors, shared array B[0..n-1] */
forall i < n {
    B[i] = A[i];
    local p = n / 2;
}
forall i < p {
    if (p == 0) stop;
    B[i] = B[2*i] + B[2*i + 1];
    p = p / 2;
} repeat
\end{lstlisting}
\end{ppprogram}

\textbf{Analysis.}
\begin{itemize}
  \item Time: $T_n = O(\log n)$ rounds.
  \item Work: $W = O(n)$ total additions (the tree has $n - 1$ internal nodes).
  \item Speedup: $S_n = n / \log n$.
  \item Efficiency: $E_n = 1 / \log n$.
  \item Cost: $C_n = n \cdot \log n$.
\end{itemize}

\begin{pppitfall}
The algorithm is \emph{not} cost-optimal: cost is $\Theta(n \log n)$ while the
sequential algorithm costs $\Theta(n)$. In later chapters we will see how to
design work-optimal parallel algorithms that achieve $O(\log n)$ time
\emph{and} $O(n)$ total work.
\end{pppitfall}

% -------------------------------------------------------------------
\section{Key Takeaways}
% -------------------------------------------------------------------

\begin{itemize}
  \item \ppterm{Speedup} $S_p = T_1/T_p$ and \ppterm{cost} $C_p = p \cdot T_p$
        are the two primary measures; a cost-optimal algorithm satisfies
        $C_p = O(T_1)$.
  \item \ppterm{Efficiency} $E_p = S_p/p$ tells you how much of each processor's
        time is doing useful work; ideal efficiency is 1. Always report it
        alongside speedup.
  \item \ppterm{Strong scaling} fixes problem size and increases $p$; it
        is bounded by the sequential fraction (Amdahl's Law) and by growing
        communication/synchronisation overheads.
  \item \ppterm{Weak scaling} grows problem size proportionally with $p$;
        ideal runtime stays constant. Gustafson's Law quantifies speedup under
        this model and shows it can grow linearly with $p$.
  \item \ppterm{Amdahl's Law}: $S_\infty = 1/f$---the serial fraction $f$
        places a hard ceiling on achievable speedup regardless of processor count.
  \item \ppterm{Gustafson's Law}: $S_p = f + (1-f)p$---as workload scales
        with $p$, the serial fraction shrinks and speedup grows linearly.
  \item The \ppterm{isoefficiency function} $W = \frac{E_c}{1-E_c} T_o$ gives
        the minimum workload growth rate needed to maintain constant efficiency;
        a lower-order isoefficiency means better scalability.
  \item \ppterm{Superlinear speedup} ($S_p > p$) is possible in practice due
        to cache effects, but cannot be predicted by the PRAM model alone.
  \item The \ppterm{PRAM model} provides a clean theoretical foundation:
        synchronous processors, constant-time shared memory, and a hierarchy of
        concurrent access models (EREW $\leq$ CREW $\leq$ CRCW).
  \item The binary-tree parallel sum achieves $O(\log n)$ time but is
        \emph{not} cost-optimal ($\Theta(n \log n)$ cost vs.\ $\Theta(n)$
        sequential).
\end{itemize}

% -------------------------------------------------------------------
\section{Exercises}
% -------------------------------------------------------------------

\begin{ppexercises}
  \ppexercise{Prove that $S_\infty = 1/f$ under Amdahl's Law. What does
  $f \to 0$ imply for scalability?}

  \ppexercise{A parallel program has $T_1 = n^2$ and $T_p = n^2/p + n \log p$.
  Compute $T_o$, find the isoefficiency function, and classify the algorithm
  as strongly or weakly scalable.}

  \ppexercise{Show that the binary-tree summation in Program~\ref{prog:pram-sum}
  requires CREW (and not just EREW). Identify the step where multiple processors
  read the same location simultaneously.}

  \ppexercise{Design an EREW variant of the parallel sum algorithm. How many
  additional steps does it require compared to the CREW version?}

  \ppexercise{Suppose you observe $S_{16} = 20$ in a benchmark. Does this
  violate Amdahl's Law? Explain what architectural phenomenon could account for
  this result and under what conditions it disappears.}

  \ppexercise{For the matrix multiplication example with $T_p = n^3/p +
  (n + 1000 \log p)$, compute the isoefficiency function assuming the desired
  efficiency is $E_c = 0.5$. How does the answer change if communication
  overhead is reduced to $100 \log p$?}
\end{ppexercises}

% -------------------------------------------------------------------
\section*{Solutions}
% -------------------------------------------------------------------

\textbf{Solution 1.}
From Amdahl's formula $S_p = 1/(f + (1-f)/p)$, take the limit as $p \to \infty$:
\[
  \lim_{p \to \infty} S_p = \lim_{p \to \infty} \frac{1}{f + (1-f)/p} = \frac{1}{f + 0} = \frac{1}{f} = S_\infty.
\]
As $f \to 0$, $S_\infty \to \infty$: there is no theoretical ceiling on speedup
when the program is fully parallelisable. In practice, communication and
synchronisation overheads prevent this limit from being reached.

\medskip
\textbf{Solution 2.}
\[
  T_o = p \cdot T_p - T_1 = p\!\left(\frac{n^2}{p} + n \log p\right) - n^2
       = n^2 + pn \log p - n^2 = pn \log p.
\]
For constant efficiency $E_c$:
\[
  W = \frac{E_c}{1-E_c} \cdot T_o = \Theta(pn \log p).
\]
Since $W = T_1 = n^2$, substituting gives $n^2 = \Theta(pn \log p)$, so
$n = \Theta(p \log p)$ and $W = \Theta(p^2 \log^2 p)$.
The workload must grow super-quadratically with $p$---the algorithm is
\ppterm{weakly scalable}, and scales poorly.

\medskip
\textbf{Solution 3.}
In Phase~1, processor $i$ writes $B[i] = A[i]$---all writes are to distinct
locations, so EREW suffices. In the reduction phase, at step $k$ processor $i$
reads both $B[2i]$ and $B[2i+1]$. At the same step, processor $i-1$ reads
$B[2(i-1)+1] = B[2i-1]$ and $B[2i]$. Hence $B[2i]$ is read simultaneously by
processors $i$ and $i-1$---this is a concurrent read. EREW does not permit
this; CREW is the minimum required model.

\medskip
\textbf{Solution 4.}
To avoid concurrent reads, use an \emph{odd-even} reduction: at step $k$,
only processor $i$ reads $B[i + 2^{k-1}]$ and adds it to $B[i]$. No two
active processors read the same location in the same step (processor $i$ reads
index $i + 2^{k-1}$, which is unique to $i$ at that step). The algorithm
still requires $\log_2 n$ steps, so the time complexity is unchanged;
the only difference is the stricter access pattern. EREW requires no extra
steps in this formulation---the same $O(\log n)$ time bound holds.

\medskip
\textbf{Solution 5.}
Amdahl's Law is a theoretical bound derived under the assumption of constant
sequential and parallel execution costs with no data-locality effects.
$S_{16} = 20 > 16$ does \emph{not} violate the \emph{law} in practice---it
violates the asymptotic model's assumptions. The mechanism is cache-induced
superlinear speedup: the full problem does not fit in a single processor's
cache, causing frequent DRAM accesses. When the problem is split across 16
processors, each sub-problem fits in its local cache, reducing effective memory
latency from $\sim$100 ns (DRAM) to $\sim$1--5 ns (L2/L3). This latency
reduction is multiplicative and not captured by the work-only model.
Superlinear speedup disappears when the single-processor data set also fits in
cache (smaller problem size), or when the parallel decomposition increases
inter-processor communication to the point where network latency outweighs
cache savings.

\medskip
\textbf{Solution 6.}
With $E_c = 0.5$: $W = \frac{0.5}{0.5} T_o = T_o$.
The overhead is $T_o = np + 1000p \log p$. For large $n$, $np$ dominates:
\[
  n^3 = \Theta(np) \implies n^2 = \Theta(p) \implies n = \Theta(\sqrt{p})
  \implies W = \Theta(p^{3/2}).
\]
(The constant factor changes from the $E_c = 0.8$ case but the asymptotic
class is the same.)

If the overhead is reduced to $T_o = np + 100p \log p$, for large $n$ the
dominant term is still $np$, so the isoefficiency remains $W = \Theta(p^{3/2})$.
The constant coefficient of the overhead shrinks, meaning the system can
maintain the target efficiency with a smaller absolute workload, but the
asymptotic growth rate is unchanged.

